%% ============================================================================
%% Elsevier CAS single-column manuscript
%% Transport-inspired rare-event sampling in spin systems
%%
%% Build: pdflatex x2 (no external .bib; thebibliography is inline)
%% Requires: cas-sc.cls (Elsevier CAS bundle)
%% ============================================================================

\documentclass[a4paper,fleqn]{cas-sc}

\usepackage[numbers,sort&compress]{natbib}
\usepackage{amsmath}
\usepackage{amssymb}
\usepackage{booktabs}
\usepackage{graphicx}
\usepackage{caption}
\usepackage{float}
\usepackage{placeins}
\usepackage{enumitem}

\setlist{itemsep=2pt,topsep=3pt}
\newcommand{\WE}{\mathrm{WE}}
\newcommand{\FOM}{\mathrm{FOM}}
% Suppress the CAS template's empty ORCID footnote; replace this with an author ORCID if available.
\RenewDocumentCommand{\printorcid}{}{}

\begin{document}
\restylefloat{figure}
\restylefloat{table}
\let\WriteBookmarks\relax
\def\floatpagepagefraction{1}
\def\textpagefraction{.001}

\shorttitle{Transport-inspired rare-event sampling in spin systems}
\shortauthors{M.Y. Alzaatreh}

\title[mode=title]{Transport-Inspired Rare-Event Sampling in Spin Systems: Learned
Importance Improves Deep-Target Reliability}

\author{M.Y. Alzaatreh}
\cormark[1]
\ead{mzaatreh@fbsu.edu.sa}
\credit{Conceptualization, Methodology, Software, Formal analysis, Investigation,
Visualization, Writing -- original draft, Writing -- review \& editing}

\affiliation{organization={Natural Sciences Unit, Fahad Bin Sultan University},
            city={Tabuk},
            country={Saudi Arabia}}

\cortext[1]{Corresponding author}

% ---- Abstract --------------------------------------------------------------
\begin{abstract}
Rare-event sampling requires computational effort to be concentrated in configurations that
disproportionately determine a target probability. Radiation transport addresses this through
adjoint importance functions and weight-window population control, and this work transfers
that toolkit to rare-event sampling in lattice spin systems, using the shared
backward-conditioning structure of the transport adjoint and the transition-path committor. A
solvable birth--death process reproduces the zero-variance limit of exact importance, and
enumeration of a $4\times4$ lattice confirms that weight-conserving Weighted Ensemble (WE)
resampling is unbiased for any binning coordinate. At matched spin-update cost and full
replication ($R=20$), WE and naive sampling resolve the $L=20$ Ising magnetization tail to
comparable depth, with WE holding a modest, roughly two-fold precision advantage at the single
deepest bin either method reaches. For the
frustrated Edwards--Anderson spin glass, a self-consistency-trained neural coordinate acting on
the full spin configuration bins the weighted population. Across six independent disorder
realizations with one trained network each, the deepest calibrated target is reached in $121$
of $180$ replicas ($67.2\%$) against $95/180$ ($52.8\%$) for the energy-binned coordinate.
A disorder-stratified Cochran--Mantel--Haenszel analysis gives a common odds ratio of $2.33$
($95\%$ CI $[1.40,3.88]$, $p=0.0011$); the same reliability advantage also appears within a
single realization across three independently trained networks compared against a matched-size
baseline ($66.7\%$ versus $48.9\%$, Fisher exact $p=0.023$). That improvement is one of
reliability rather than wall-clock efficiency. Repeated neural inference costs four to five
times the wall-clock time of the energy coordinate, so the directly measured figure of merit
sits at roughly one-third of parity ($0.33$, $95\%$ CI $[0.19,0.58]$). These results support
the transfer of the tested transport-inspired framework and show that the reliability advantage
extends beyond a single quenched-disorder realization. At the tested scale, neural inference
overhead is the dominant measured barrier to wall-clock competitiveness.
\end{abstract}

% ---- Highlights ------------------------------------------------------------
\begin{highlights}
\item Adjoint importance and the rare-event committor share the same backward-conditioning
structure and coincide for the unit-response absorbing-target problem studied here.
\item Exact tests validate both the zero-variance importance limit and unbiased WE population
control.
\item At matched update cost and full replication ($R=20$), WE holds an approximately
two-fold precision advantage over naive sampling at the deepest resolved Ising-tail bin,
though not uniformly across the tail.
\item Learned importance improves deep-event discovery across six independent
Edwards--Anderson disorder realizations ($67.2\%$ versus $52.8\%$; stratified odds ratio $2.33$,
$p=0.0011$).
\item The reliability improvement does not yet translate into wall-clock efficiency
($0.33$, $95\%$ CI $[0.19,0.58]$).
\end{highlights}

% ---- Keywords --------------------------------------------------------------
\begin{keywords}
Rare-event sampling \sep committor \sep adjoint importance \sep Weighted Ensemble \sep
variance reduction \sep Ising model \sep spin glass \sep neural importance
\end{keywords}

\maketitle

% ============================================================================
\section{Introduction}\label{sec:intro}

A rare event is one that direct simulation almost never observes: the nucleation of a new
phase, a protein crossing between conformational basins \citep{bolhuis2002}, or a detector
registering a particle that has penetrated many mean free paths of shielding
\citep{kahn1951,kalos2008}. In each case the probability
of the event is itself the quantity of scientific interest, and it is exactly that quantity
which unbiased dynamics represents worst. If an event has equilibrium or first-passage
probability $p\ll1$, the relative error of a direct Bernoulli estimator formed from $N$
effectively independent samples follows the elementary result \citep{kahn1951,kalos2008}
\begin{equation}
\varepsilon_{\rm rel}=\sqrt{\frac{1-p}{Np}}\;\simeq\;(Np)^{-1/2},\qquad p\ll1,
\label{eq:bernoulli}
\end{equation}
so the sample count required for fixed relative precision grows as $1/p$ and each additional
decade of rarity costs an order of magnitude in computation. The failure is structural rather
than incidental, and every established remedy responds to it the same way: not with more
samples, but with better-placed ones. This elementary scaling is the origin of a broad family
of rare-event techniques, including importance sampling,
splitting, transition-path sampling, Forward Flux Sampling (FFS), Adaptive Multilevel
Splitting (AMS), and Weighted Ensemble (WE) methods \citep{kahn1951,kalos2008,huber1996,
zuckerman2017,allen2009,cerou2007,bolhuis2002}. What each of them requires in practice is a
rule for deciding which samples are the better-placed ones.

Two mature traditions answer that question using closely related mathematical objects. In
transition-path theory, the committor $h(s)$ is the probability that a trajectory initiated
at configuration $s$ reaches a target set before a competing boundary, and it provides an
optimal reaction coordinate for that event \citep{e2006,metzner2009,vandeneijnden2010}. In
classical particle transport, an adjoint solution assigns an importance to a phase-space
point according to its expected contribution to a detector response; this importance can be
used to bias sources, alter transition probabilities, and construct weight windows
\citep{lewins1965,spanier1969,bell1970,wagner1998,haghighat2003}. The formal connection to
the Doob $h$-transform and conditioned Markov processes is well known
\citep{doob1984,touchette2009,chetrite2015}. The contribution of the present work is not to
claim a new identity between these objects, but to carry out a systematic transfer and
validation of the transport variance-reduction viewpoint in spin-system rare-event sampling,
together with a controlled test of a learned configuration-space coordinate built on this
viewpoint.

The transport perspective is attractive because it separates the ideal importance function
from the population-control mechanism used when that importance is only approximate. Exact
adjoint biasing can in principle produce zero-variance estimators, while weight windows use
splitting and roulette to keep statistical weight near a prescribed range. This distinction
underlies CADIS and related hybrid deterministic--Monte Carlo variance-reduction methods
\citep{wagner1998,haghighat2003,wagnerfw2014}. In the present setting, the same logic
motivates using a committor-like quantity to organize a weighted population while retaining
the original spin dynamics.

Rare-event population methods have a parallel history outside transport. WE was introduced
for biomolecular transition sampling by Huber and Kim, shown to be statistically exact for
broad classes of dynamics and binning procedures, and made practical for complex molecular
systems by steady-state formulations and mature software
\citep{huber1996,zhang2010,bhatt2010,russo2022}; FFS and AMS address related first-passage
problems through interfaces or adaptive levels \citep{allen2009,valeriani2007,cerou2007},
while transition-path sampling works with reactive-path ensembles and can avoid prescribing a
one-dimensional coordinate at all \citep{dellago1998,bolhuis2002}. Machine learning
meanwhile offers several routes to the high-dimensional objects these methods need: flow-based
and autoregressive samplers for equilibrium distributions
\citep{wu2019,noe2019,nicoli2020,gabrie2022}, neural approximations of committor and
Feynman--Kac observables from short trajectories \citep{khoo2019,strahan2023}, learned bias
potentials and importance functions \citep{hua2024,kimcai2026}, and neural Doob dynamics for
large deviations \citep{zhao2025}. Together these motivate the central practical question of
this paper: when the rare event occurs in a high-dimensional spin configuration space, how
much is gained by replacing a hand-designed coordinate with a learned, committor-inspired one,
measured against a strong physical baseline and with inference cost charged to the result?

The study is organized as a validation ladder rather than a single benchmark, built around
three questions. Is the transport/committor correspondence mathematically valid in the
formulated first-passage setting? How does weight-conserving WE alter sampling precision
relative to direct Monte Carlo as the Ising tail becomes progressively rarer? And does a learned
full-configuration coordinate improve deep-event access over a strong physical energy
coordinate, reproducibly across independent instances of the problem rather than for one
particular realization? A solvable
one-dimensional birth--death chain first verifies the zero-variance limit of exact importance.
Exact enumeration of a $4\times4$ Ising lattice then verifies unbiased WE resampling. A larger
$L=20$ Ising calculation tests the rarity-dependent precision trade-off between WE and direct
sampling at matched spin-update cost, and finally the Edwards--Anderson (EA) spin glass provides a
frustrated landscape in which magnetization is a poor coordinate and the low-energy target is
a more demanding test for learned configuration-space importance
\citep{ising1925,onsager1944,edwards1975,binder1986}. Throughout, \emph{reliability} (whether
a method reaches the event at all, reproducibly) is distinguished from \emph{efficiency}
(variance reduction per unit of actual computational cost). That distinction matters most for
neural coordinates, since inference overhead is not represented by a spin-update count.

The main findings are positive but bounded. WE is unbiased, validated against exact
references, and holds a modest precision advantage over naive sampling at the deepest
resolved point of the Ising tail, though not uniformly across it. In the EA model, a
self-consistency-trained neural coordinate significantly improves deep-event discovery over
the physical energy coordinate, an advantage reproduced both across independently trained
networks on one disorder realization and across six independent disorder realizations. That
improvement is a reliability result, though: once inference is included in the cost, the same
coordinate does not improve the figure of merit. Together, these results identify a
transferable theoretical core, a validated route for using learned importance information
inside an unbiased weighted-population sampler, while identifying neural inference overhead as
the dominant measured barrier to wall-clock competitiveness at the tested scale.

\FloatBarrier

% ============================================================================
\section{Theory and Methods}\label{sec:methods}

\subsection{Adjoint importance, committors, and zero-variance sampling}\label{sec:adjoint}
Let $K(s,s')$ be the transition kernel of a Markov chain and let $B$ denote a target set.
For a first-passage event, the forward committor is the probability of reaching $B$ before
the competing set $A$ \citep{e2006,metzner2009}:
\begin{equation}
h(s)=\Pr_{s}\!\left(\tau_B<\tau_A\right).
\label{eq:committor}
\end{equation}
For a discrete Markov process, this committor satisfies the backward (harmonic) equation in
the interior of the state space, a standard result of transition-path theory for Markov jump
processes \citep{e2006,metzner2009,vandeneijnden2010}:
\begin{equation}
h(s)=\sum_{s'}K(s,s')\,h(s'),\qquad s\notin A\cup B,
\label{eq:harmonic}
\end{equation}
with boundary values $h=0$ on $A$ and $h=1$ on $B$. Conditioning the dynamics with a
positive harmonic function leads to the Doob $h$-transform, introduced in classical potential
theory \citep{doob1984} and used in this form throughout the large-deviation and
conditioned-process literature \citep{touchette2009,chetrite2015}:
\begin{equation}
K_h(s,s')=K(s,s')\,\frac{h(s')}{h(s)},
\label{eq:doob}
\end{equation}
where defined. This is the Markov-chain counterpart of adjoint-based importance biasing in
transport. With the exact $h$, successful paths acquire identical total statistical weight,
producing the familiar zero-variance ideal.

The relationship between the two importance functions should be stated precisely, since it is
weaker than an identity. The committor of Eq.~\eqref{eq:committor} is a probability: the chance
of reaching $B$ before $A$ from $s$, bounded in $[0,1]$ and defined only once a competing
boundary is specified. The transport adjoint, by contrast, is an expected contribution to a
specified detector response; it carries the units of that response and is unbounded in
general. What the two share is structure rather than definition. Both are solutions of a
backward equation generated by the same transition operator, both express the value of a state
with respect to a terminal objective, and both play the same role when used to tilt or to bin.
They coincide as functions in the particular case of an absorbing target with unit detector
cross-section and no multiplication, which is the case constructed here, but they are not
universally the same object and this paper does not claim they are. What transfers is the
backward-conditioning structure and the machinery built on it, not an equivalence of the two
quantities in general.

This identity motivates two distinct uses of approximate importance. The first changes
transition probabilities directly, which can be fragile when $h$ is inaccurate. The second
retains the original dynamics but uses population splitting and merging to control statistical
weight. The second strategy is adopted here using an unbiased, weight-conserving WE resampling
scheme, so an imperfect binning coordinate can alter variance and sampling allocation without
altering the estimator expectation \citep{zhang2010}.

The exact equivalence and zero-variance construction above apply to the formulated first-passage
problem. In the equilibrium tail calculations below, the committor-inspired quantity is not
used to claim an exact zero-variance transformation of the stationary distribution; it is used
only as a progress coordinate for unbiased WE resampling. The first-passage theory therefore
supplies the importance principle, while the spin-system calculations test its utility as a
population-allocation coordinate.

\subsection{Zero-variance path weighting}\label{sec:zerovar}
The zero-variance statement can be seen directly at the path level. Consider a trajectory
$\omega=(s_0,s_1,\dots,s_\tau)$ generated under the Doob-transformed kernel until it reaches
either $A$ or $B$. Substituting the Doob transform of Eq.~\eqref{eq:doob} into the path
likelihood ratio gives the telescoping identity below, derived here for the discrete
first-passage setting from Eqs.~\eqref{eq:harmonic}--\eqref{eq:doob}; it is the discrete-state
counterpart of the classical zero-variance principle of transport importance sampling
\citep{lewins1965,spanier1969,haghighat2003} and of the conditioned-process construction of
Doob \citep{doob1984} and Chetrite and Touchette \citep{chetrite2015}:
\begin{equation}
\prod_{t=0}^{\tau-1}\frac{K(s_t,s_{t+1})}{K_h(s_t,s_{t+1})}
=\prod_{t=0}^{\tau-1}\frac{h(s_t)}{h(s_{t+1})}
=\frac{h(s_0)}{h(s_\tau)} .
\label{eq:telescope}
\end{equation}
On a successful path ending in $B$, $h(s_\tau)=1$, so every successful history contributes
exactly $h(s_0)$.

Fix an initial state $s_0$ with $h(s_0)>0$ and let the target quantity be the first-passage
probability $\pi=h(s_0)=\Pr_{s_0}(\tau_B<\tau_A)$. Generate $M$ independent trajectories under
$K_h$ and assign each the likelihood-ratio weight of Eq.~\eqref{eq:telescope}, which gives
the following estimator, derived here:
\begin{equation}
\hat\pi=\frac1M\sum_{m=1}^{M}
\left[\prod_{t=0}^{\tau_m-1}\frac{K(s_t^{(m)},s_{t+1}^{(m)})}{K_h(s_t^{(m)},s_{t+1}^{(m)})}\right]
\mathbf 1_B\!\left(s_{\tau_m}^{(m)}\right).
\label{eq:zvest}
\end{equation}
Here $\mathbf 1_B(\cdot)$ denotes the indicator function of the target set $B$, equal to $1$ if
its argument lies in $B$ and $0$ otherwise.
States with $h=0$ receive zero transition probability, so the transformed chain never enters
$A$ and every trajectory terminates in $B$; on each such trajectory the weight telescopes to
$h(s_0)$. Every term of the sum therefore equals $h(s_0)$, so $\hat\pi=\pi$ identically and
$\mathrm{Var}(\hat\pi)=0$ for any $M\ge1$. This holds for the exact $h$; it degrades
continuously, and in general non-uniformly, when $h$ is replaced by an approximation. This is
the property verified numerically in \S\ref{sec:bd}.

For an approximate importance function $\tilde h$, the exact telescoping property is no
longer available. The present implementation separates importance representation from
dynamical biasing instead: the physical Metropolis kernel is retained, while the learned or
physical progress coordinate is used only to organize population splitting and merging. This
design lets approximate importance information guide computational effort while weight
conservation preserves the target measure.

Table~\ref{tab:map} summarizes the operational correspondence used throughout the paper. The
mapping is not intended to imply that all transport constructs carry over unchanged; source
biasing, for example, has no literal analogue when the equilibrium spin distribution is fixed
and the goal is to sample a stationary tail. The useful transfer is structural: an
adjoint-like scalar identifies configurations that matter to the response, and a
weight-management mechanism concentrates trajectory population in those configurations.

\begin{table}[H]
\centering\small
\caption{Transport concepts and their rare-event spin-system counterparts used in this work.}
\label{tab:map}
\begin{tabular}{p{0.26\textwidth}p{0.28\textwidth}p{0.34\textwidth}}
\toprule
Radiation-transport concept & Markov/spin counterpart & Role in this study \\
\midrule
Adjoint importance & Committor / event-conditioned value function & Defines ideal progress
toward a rare target \\
Zero-variance importance bias & Doob $h$-transform / conditioned dynamics & Verified in the
solvable first-passage demonstration \\
CADIS / weight-window construction & Importance-guided population allocation & Motivates the
use of an importance coordinate to organize the weighted population \\
Weight windows & Weighted splitting and merging & Implemented through unbiased WE resampling
with conserved statistical weight \\
Detector response & Rare-event indicator or threshold probability & Magnetization-tail or
low-energy probability \\
\bottomrule
\end{tabular}
\end{table}

\subsection{Spin models and rare-event targets}\label{sec:models}
For the nearest-neighbor ferromagnetic Ising model, the Hamiltonian is that introduced by
Ising \citep{ising1925} and solved exactly in two dimensions by Onsager \citep{onsager1944},
\begin{equation}
H(\mathbf s)=-J\sum_{\langle ij\rangle}s_is_j,\qquad s_i\in\{-1,+1\},
\label{eq:ising}
\end{equation}
with periodic boundaries and $J=1$. The principal rare observable is the absolute
magnetization $|m|$. For the spin glass, the corresponding random-bond Hamiltonian is that of
Edwards and Anderson \citep{edwards1975}, in the $\pm J$ form reviewed by Binder and Young
\citep{binder1986},
\begin{equation}
H(\mathbf s)=-\sum_{\langle ij\rangle}J_{ij}s_is_j,\qquad J_{ij}\in\{-1,+1\},
\label{eq:ea}
\end{equation}
with one fixed disorder realization per comparison (six independent realizations are
compared in \S\ref{sec:disordergen}). Quenched disorder and frustration make the low-energy
states disordered, so $m\approx0$ irrespective of energy and magnetization is not an adequate
progress coordinate; low energy is used as the rare-event target instead. All reported
production comparisons use single-spin Metropolis dynamics \citep{metropolis1953} at $T=2.6$
unless stated otherwise.

\subsection{Weighted Ensemble estimator and resampling}\label{sec:we}
WE propagates a weighted population under the original, unbiased Metropolis dynamics for a
resampling interval $\tau$. After propagation, walkers are partitioned into bins defined by a
coordinate $q(s)$. If the total statistical weight in bin $b$ is $W_b$, resampling
redistributes that same $W_b$ among the target number of walkers in the bin. Splitting creates
additional trajectories with reduced individual weights this way, while merging removes
redundant trajectories without changing total weight. The split/merge rule used here is an unbiased, weight-conserving WE resampling scheme
\citep{zhang2010}. Consequently, the coordinate controls where sampling effort is placed and can
change estimator variance, while the expectation of the target estimator is preserved.

For a target set $B$, WE estimates the target probability from the total statistical weight
carried by walkers in $B$. The estimator takes the weight-conserving form established for WE
by Huber and Kim \citep{huber1996}, proved statistically exact for broad classes of dynamics
and binning procedures by Zhang et al.\ \citep{zhang2010}, and used in the steady-state and
review formulations of Bhatt et al.\ \citep{bhatt2010} and Zuckerman and Chong
\citep{zuckerman2017}:
\begin{equation}
\hat\pi_t=\frac{\sum_i w_i^{(t)}\,\mathbf 1_B\!\left(s_i^{(t)}\right)}{\sum_i w_i^{(t)}},
\label{eq:weest}
\end{equation}
Here $i$ indexes the walkers in the population at resampling cycle $t$, $w_i^{(t)}$ and
$s_i^{(t)}$ are the statistical weight and state of walker $i$, and $\mathbf 1_B(\cdot)$ is the
same indicator function of the target set $B$ used in Eq.~\eqref{eq:zvest}.
The reported estimate is averaged over post-burn-in resampling cycles. Three coordinates
are compared: absolute magnetization $|m|$, energy per spin $E/N$, and a learned scalar
coordinate $I_\theta(s)$. Magnetization is deliberately retained as a mismatched baseline for
the EA spin glass, whereas energy is a natural physical baseline for the low-energy target.

A fixed rare-event threshold can be either trivial or physically unreachable when lattice
size, temperature, or disorder changes. The computational protocol uses a short pilot
simulation instead, to locate the most extreme value reached under ordinary dynamics, and then
places the production threshold beyond that pilot extreme by a small number of discrete
lattice steps (spacing $2/L^2$ for magnetization, $4/L^2$ for EA energy per spin). The integer
``beyond'' records how many such steps the target is displaced past the pilot extreme; it
provides an operational depth scale, without claiming that equal values of ``beyond''
correspond to identical probabilities across different system sizes or disorder realizations.
The resampling interval is treated as a physical parameter rather than a bookkeeping choice.
Independent diagnostics in the EA model show rapid relaxation of deep configurations, so
shorter resampling intervals penetrate the deep ladder more effectively at comparable cost;
the production deep-rare calculation uses $\tau=2$ for this reason.

\subsection{Neural importance coordinate}\label{sec:neural}
The learned coordinate is a compact periodic convolutional network acting on the full spin
configuration: two circularly padded $3\times3$ convolutional layers with nonlinear
activations, followed by global mean and maximum pooling and a scalar output. This
architecture is intentionally small: the research question is whether a configuration-space
coordinate can discover useful progress information absent from a simple scalar observable,
not whether a large neural model can fit the training set.

The training objective is shaped by the harmonic self-consistency condition of
Eq.~\eqref{eq:harmonic} itself, in the spirit of the recent label-free neural
importance-sampling formulation of Kim and Cai \citep{kimcai2026}, in which a learned
importance function is coupled to biased transitions and branching. This direction also sits
within a broader neural-committor literature: direct variational approximation of
high-dimensional committors, semigroup formulations avoiding explicit second derivatives,
active importance sampling for rare-event-dominated objectives, short-trajectory
Feynman--Kac training, and self-consistent committor-enhanced sampling
\citep{khoo2019,li2022,rotskoff2022,strahan2023,kang2024,trizio2025}. The loss penalizes
deviations from the harmonic self-consistency relation and enforces both a target boundary and
a distinct bulk-reference boundary; the latter is required because the harmonic condition plus
the target boundary alone admits a trivial constant solution. Production training subsets were
therefore constructed to retain both boundary-adjacent and bulk-reference configurations, with
boundary-adjacent states capped at $70\%$ of the $n_{\rm train}=2500$ subset.

For the full production preset, each convolutional layer uses $24$ channels and ReLU
activations. Training uses Adam with learning rate $2\times10^{-3}$, mini-batches of $128$, and
$20$ epochs. The network output is used only as the WE binning coordinate. This is more
conservative than the Kim--Cai construction: coordinate quality can be evaluated without
placing an approximate learned function inside a likelihood-ratio tilt.

The trained network $I_\theta$ is accordingly termed a \emph{learned importance coordinate},
not a neural committor: its objective is shaped by the harmonic condition of
Eq.~\eqref{eq:harmonic}, but it is not claimed to be a quantitatively accurate estimate of $h$,
and no such accuracy is measured here. Because $I_\theta$ enters only through the binning coordinate and the WE resampling rule is
unbiased for arbitrary bins under the conditions established for WE \citep{zhang2010}, the
results below do not require $I_\theta$ to be quantitatively close to $h$: an inaccurate
coordinate can degrade efficiency without changing the estimator expectation.

\subsection{Computational cost, statistics, and reporting}\label{sec:cost}
Two cost measures are used. Matched-cost comparisons of sampling reach use the number of
attempted spin updates as the common Monte Carlo work unit, natural for the Ising-tail
comparison since both direct sampling and WE use the same Metropolis kernel. That measure is
blind to neural inference, though, which is performed repeatedly during WE binning; wall-clock
time is also recorded directly for this reason, and is the measure used whenever a figure of
merit is quoted for a neural coordinate. The figure of merit is the cost-normalized inverse
variance standard in transport variance-reduction studies \citep{spanier1969,wagnerfw2014},
\begin{equation}
\FOM=\frac{1}{\sigma_{\rm rel}^{2}\,T},
\label{eq:fom}
\end{equation}
with $T$ the computational cost charged to the calculation.

Suppose $R$ independent replicas produce estimates $\hat\pi_r$. Following the
independent-replica uncertainty reporting recommended for WE studies \citep{zuckerman2017} and
implemented in WESTPA \citep{russo2022}, the across-replica relative standard deviation is
\begin{equation}
\sigma_{\rm rel}=\frac{\mathrm{sd}(\hat\pi_1,\dots,\hat\pi_R)}{\bar\pi},\qquad
\bar\pi=\frac1R\sum_{r=1}^{R}\hat\pi_r ,
\label{eq:relsd}
\end{equation}
but for deep targets the most informative quantity can instead be the fraction of independent
runs that observe the event at all, so zero-hit counts are reported explicitly. Event-positive
and zero-hit replicas are compared using the exact conditional test for a $2\times2$
contingency table introduced by Fisher \citep{fisher1922}, appropriate here because the
expected counts are small. For the six-disorder comparison, the descriptive pooled Fisher test is
complemented by a Cochran--Mantel--Haenszel analysis that preserves disorder realization as a
stratum and estimates a common odds ratio \citep{mantel1959}. Confidence intervals on
figure-of-merit ratios are obtained by the nonparametric bootstrap \citep{efron1979} in a
two-level form that resamples the trained-network (or disorder-realization) axis in the outer
loop and the replica axis in the inner loop.

During software development, ChatGPT (OpenAI) and Claude (Anthropic) were used as coding
assistants for code drafting, refactoring, and diagnostic review. Production calculations were
executed under the recorded configurations, and code changes affecting reported results were
checked against the exact validation gates described in \S\ref{sec:validation}.

\FloatBarrier

% ============================================================================
\section{Validation}\label{sec:validation}

\subsection{Exact-importance limit}\label{sec:bd}
A biased one-dimensional birth--death chain provides a solvable test. Starting at $s_0=1$
with absorbing boundaries at $0$ and $N=15$ and step probability $p=0.4$ toward the rare
boundary, the exact committor is $h(s_0)=1.144443\times10^{-3}$. Direct Monte Carlo gives a
relative variance of $2.17\times10^{-3}$ for $4\times10^{5}$ histories, in agreement with the
Bernoulli prediction $1/(Np)$. Using the exact importance transform reduces the measured
relative variance to approximately $10^{-30}$, numerically reproducing the zero-variance
prediction of Eq.~\eqref{eq:telescope}.

\subsection{Exact-enumeration gate for WE}\label{sec:gate}
For the $4\times4$ Ising model, all $2^{16}$ configurations can be enumerated. Across three
validation regimes (above $T_c$, below $T_c$ with the double-well structure, and the folded
$|m|$ distribution), the mean $|\Delta\ln P|$ between WE and exact enumeration ranges from
approximately $0.006$ to $0.024$; a separate learned-importance gate gives target-probability
ratios of $1.000$ and $1.001$ relative to exact enumeration when WE is binned by magnetization
and energy, respectively. These tests confirm that the implemented WE resampling reproduces the exact target measure
under distinct binning coordinates, so later coordinate comparisons are interpreted as
efficiency or reliability differences rather than changes in the target expectation.

\begin{table}[H]
\centering\small
\caption{Validation ladder used before interpreting the larger rare-event calculations.}
\label{tab:ladder}
\begin{tabular}{p{0.30\textwidth}p{0.24\textwidth}p{0.34\textwidth}}
\toprule
Test & Reference & Outcome \\
\midrule
Birth--death importance transform & analytic committor & exact-importance relative variance
$\sim10^{-30}$ \\
$4\times4$ Ising WE distribution & exact enumeration & mean $|\Delta\ln P|=0.006$--$0.024$
across three regimes \\
$4\times4$ learned-coordinate gate & exact target probability & WE ratios $1.000$ and
$1.001$ \\
\bottomrule
\end{tabular}
\end{table}

\FloatBarrier

% ============================================================================
\section{Rare-Tail Sampling in the Ising Ferromagnet}\label{sec:tail}

\subsection{WE and naive sampling in the Ising tail}\label{sec:tailreach}
At $L=20$ and matched update cost ($\sim3.4\times10^{8}$ spin updates), naive sampling and WE
agree closely through the bulk and into the near tail, with $R=20$ independent replicas per
method. At the second-deepest resolved bin ($|m|=0.886$), naive sampling is in fact the more
precise of the two at matched cost: relative standard deviation $0.36$ against WE's $0.79$, a
figure-of-merit advantage for naive of roughly $5\times$. Only at the single deepest bin
reached by either method ($|m|=0.932$, $P\approx1.1\times10^{-7}$ for WE against
$P\approx7.4\times10^{-8}$ for naive) does WE hold an edge, with a figure of merit about
$2\times$ that of naive (relative standard deviation $3.10$ versus $4.47$); both estimates are
individually noisy at this depth.

The comparison therefore shows a rarity-dependent crossover in precision rather than a
uniform algorithmic advantage. Direct sampling is more efficient in the near tail, where the
event is still observed often enough that population control adds overhead, whereas WE becomes
more precise at the deepest resolved bin, where direct-sampling variance has grown sharply.
The exact gates of \S\ref{sec:validation}, rather than this performance comparison, establish
the correctness of the WE implementation.

\subsection{Calibration of the deep-rare spin-glass regime}\label{sec:depth}
The neural-coordinate experiment is placed deliberately beyond the range where direct sampling
is consistently informative. Table~\ref{tab:depth} summarizes the depth calibration at $L=16$
and identifies the access-dominated regime used for the learned-coordinate test.

\begin{table}[H]
\centering\small
\caption{Depth calibration for the EA model at $L=16$ using 10 independent direct-sampling
replicas. The integer ``beyond'' indexes successively deeper discrete energy thresholds.}
\label{tab:depth}
\begin{tabular}{lrrr}
\toprule
Beyond & $\hat\pi$ & relative s.d. & zero-hit replicas \\
\midrule
1 & $5.29\times10^{-5}$ & 0.23 & 0/10 \\
2 & $5.48\times10^{-6}$ & 0.46 & 0/10 \\
3 & $7.14\times10^{-7}$ & 1.61 & 7/10 \\
4 & $0$ & --- & 10/10 \\
\bottomrule
\end{tabular}
\end{table}

The beyond$=4$ threshold is used as a stress test of rare-event access for this reason: direct
Monte Carlo is fully blind there, so the learned-coordinate comparison directly tests whether a
configuration-space importance map can organize the weighted population so that trajectories
continue to enter a region essentially absent from the direct Monte Carlo sample.

\FloatBarrier

% ============================================================================
\section{Learned Importance in the Deep-Rare Spin Glass}\label{sec:deep}

At beyond$=4$, direct Monte Carlo is effectively blind under the calibrated budget. The WE
budget was increased to 80 bins, 40 walkers per occupied bin, 3000 iterations, and $\tau=2$.

\subsection{One disorder, multiple trained networks}\label{sec:deepwithin}
Three independently trained self-consistency networks were compared with the same
physical-coordinate baselines on one EA disorder realization ($L=16$, $T=2.6$), $30$ replicas
per method. Using independently trained networks establishes a reproducible coordinate-level
effect rather than a favorable outcome from one particular initialization.

\begin{table}[H]
\centering\small
\caption{Deep-rare EA comparison at $L=16$, beyond$=4$, one disorder realization. Each
learned-coordinate row is a separately trained network evaluated over 30 replicas.}
\label{tab:deep}
\begin{tabular}{lrr}
\toprule
Method & event-positive replicas & discovery rate \\
\midrule
naive & 1/30 & 3.3\% \\
$\WE[m]$ & 5/30 & 16.7\% \\
$\WE[E]$ (balanced, $n=90$) & 44/90 & 48.9\% \\
$\WE[I_\theta]$, seed 0 & 19/30 & 63.3\% \\
$\WE[I_\theta]$, seed 1 & 20/30 & 66.7\% \\
$\WE[I_\theta]$, seed 2 & 21/30 & 70.0\% \\
$\WE[I_\theta]$, pooled ($n=90$) & 60/90 & 66.7\% \\
\bottomrule
\end{tabular}
\end{table}

The three independently trained networks give discovery rates of $63.3\%$, $66.7\%$, and
$70.0\%$. An expanded independent $\WE[E]$ baseline of $90$ replicas is used for the balanced
comparison, giving $44/90$ ($48.9\%$) against the pooled learned-coordinate result of $60/90$
($66.7\%$). Fisher's exact test on this balanced comparison gives
\begin{equation}
\WE[I_\theta]:\ 60/90\ (66.7\%)\quad\text{versus}\quad\WE[E]:\ 44/90\ (48.9\%),\qquad p=0.023.
\label{eq:withinp}
\end{equation}
This is the central within-disorder result: the learned importance map significantly increases
the probability that a finite-budget WE run reaches the deep target, reproducibly across three
independent initializations.

The comparison is informative because energy is a strong physical baseline for this task: the
target itself is defined by low energy, so an improvement over $\WE[E]$ is not a comparison
against an irrelevant coordinate. The network receives the full spin configuration and can
distinguish configurations with similar instantaneous energy but different local structure and
different dynamical prospects of progressing toward the target. Two configurations at
identical energy can have very different committor values on a frustrated landscape, after
all, and an energy-adjacent label cannot represent that difference. The role of the learned
function remains conservative throughout: it defines the WE partition rather than replacing
the physical transition kernel, so the learned model influences how computational effort is
distributed across configuration space while the unbiased WE resampling rule preserves the
target expectation.

\subsection{Generalization across disorder realizations}\label{sec:disordergen}
The result above establishes reproducibility across network initializations on one $J_{ij}$
sample; it does not by itself show that the effect extends beyond that particular disorder
realization. Five further independent EA disorder realizations were therefore tested at the same
$L=16$, $T=2.6$, beyond$=4$ configuration, one trained network each. Because the
within-disorder experiment already probes initialization robustness, the additional computational
budget was allocated to disorder variation rather than to another three-network ensemble for
each realization.

\begin{table}[H]
\centering\small
\caption{Deep-rare EA comparison repeated across six independent disorder realizations, one
trained network each, beyond$=4$, $30$ replicas per method. Realization 12345 is the
single-network reading of Table~\ref{tab:deep} (seed~0), not the three-network pooled result.}
\label{tab:disorder}
\begin{tabular}{lrrr}
\toprule
Disorder realization & $\WE[I_\theta]$ & $\WE[E]$ & Fisher $p$ \\
\midrule
12345 (original) & 19/30 (63.3\%) & 12/30 (40.0\%) & 0.120 \\
23456 & 15/30 (50.0\%) & 8/30 (26.7\%) & 0.110 \\
45678 & 29/30 (96.7\%) & 26/30 (86.7\%) & 0.353 \\
56789 & 20/30 (66.7\%) & 15/30 (50.0\%) & 0.295 \\
67890 & 28/30 (93.3\%) & 29/30 (96.7\%) & 1.000 \\
78901 & 10/30 (33.3\%) & 5/30 (16.7\%) & 0.233 \\
\midrule
Pooled (6 realizations) & 121/180 (67.2\%) & 95/180 (52.8\%) & 0.0071 \\
\bottomrule
\end{tabular}
\end{table}

\begin{figure}[H]
\centering
\includegraphics[width=\linewidth]{figs/fig_deeprare.pdf}
\caption{The learned importance coordinate reaches the deep target more reliably than the
physical energy coordinate. Error bars are Wilson $95\%$ intervals. (a) Discovery rate at
beyond$=4$ (disorder realization $12345$): open diamonds are the three independently trained
networks (Table~\ref{tab:deep}); $\WE[E]$ uses the balanced $90$-replica baseline. (b) The same
comparison repeated across six independent disorder realizations, one trained network each
(Table~\ref{tab:disorder}), with the pooled result at right.}
\label{fig:deeprare}
\end{figure}

No single realization reaches significance on its own at $n=30$. Realizations 45678 and 67890
are near ceiling for both methods, which limits their discriminating power; realization 67890
is also the only stratum in which $\WE[E]$ has the nominally higher discovery rate ($96.7\%$
versus $93.3\%$). In the other five of six realizations the learned coordinate has the higher
discovery rate. For descriptive comparison, pooling all six realizations gives:
\begin{equation}
\WE[I_\theta]:\ 121/180\ (67.2\%) \quad\text{versus}\quad \WE[E]:\ 95/180\ (52.8\%),
\qquad p=0.0071.
\label{eq:disordergen}
\end{equation}
Because the six comparisons have different disorder-specific event rates, inferential analysis
also preserves disorder as a stratum. The Cochran--Mantel--Haenszel analysis gives a common odds
ratio of $2.33$ ($95\%$ CI $[1.40,3.88]$) with $p=0.0011$ \citep{mantel1959}, corroborating the
descriptive pooled result without treating the six disorder realizations as one homogeneous
sample. Together with the within-disorder network replication, this shows that the reliability
advantage extends beyond a single quenched-disorder realization and is reproduced across the
six realizations tested. Wall-clock cost ratios for the five additional realizations
($4.45$--$4.97$ relative to $\WE[E]$) are consistent with the other measurements in this study
(\S\ref{sec:costdisc}), so the cost penalty is likewise reproduced across disorder.

\FloatBarrier

% ============================================================================
\section{Discussion}\label{sec:discussion}

\subsection{What is established}\label{sec:established}
Five elements of the transport-inspired framework are established by direct test. Exact
importance gives the zero-variance first-passage ideal (\S\ref{sec:bd}). Weight-conserving WE
is unbiased in the exact enumeration and learned-coordinate gates (\S\ref{sec:gate}). WE holds
a real, roughly two-fold precision advantage over naive sampling at the deepest resolved point
of the Ising tail at matched update cost (\S\ref{sec:tailreach}). The learned importance
coordinate improves deep-target discovery relative to energy in the EA spin glass
(\S\ref{sec:deepwithin}). That reliability advantage is reproduced across six independent
disorder realizations and is not confined to the original quenched-disorder sample
(\S\ref{sec:disordergen}).

\subsection{What is not established}\label{sec:notestablished}
The learned coordinate is not wall-clock superior (\S\ref{sec:costdisc}). WE's precision
advantage over naive sampling in the Ising tail is confined to the single deepest resolved bin,
not a uniform improvement across the tail (\S\ref{sec:tailreach}). Only one lattice size, one temperature, and a modest ($n=6$) disorder
ensemble have been studied for the main learned-coordinate result, and that ensemble is not
large enough to characterize how the effect's magnitude varies with disorder. The learned
coordinate has not been demonstrated to be quantitatively close to the true committor; only
that its use as a binning coordinate improves access, which does not require that accuracy
(\S\ref{sec:neural}).

\subsection{Why this matters}\label{sec:whymatters}
A coordinate can be dynamically more informative, in the sense of organizing a weighted
population so that it reaches a target more often, while being computationally less efficient
once every real cost is counted. Reliability and efficiency are not the same axis. Conflating
them would have overstated the practical readiness of the learned coordinate at this scale
while understating its scientific content: on a frustrated, quenched-disordered landscape
where the scalar energy coordinate is physically informative but compresses the configuration
to one number, the learned coordinate demonstrably extracts additional progress information.

\subsection{Reliability is improved; cost-normalized efficiency is not}\label{sec:costdisc}
On the spin-update cost measure, the three within-disorder trained networks give
figure-of-merit ratios of $1.29$, $1.07$, and $1.04$ against $\WE[E]$, bootstrapping to a
geometric-mean ratio of $1.16$ ($95\%$ CI $[0.65,2.12]$) that does not exclude parity. That
measure does not charge for network evaluations performed at every resampling step, though, so
wall-clock time was recorded directly (not reconstructed from log timestamps) for every
method. Measured ratios of $\WE[I_\theta]$ to $\WE[E]$ wall-clock time were $4.65$, $4.01$, and
$4.54$ for the three within-disorder networks, and $4.45$, $4.97$, $4.93$, $4.80$, and $4.48$
for the five additional disorder realizations: consistently four to five times $\WE[E]$'s
cost. A two-level bootstrap
gives the wall-clock-fair figure below, resampling the native $30$-replica count from the
balanced $90$-replica $\WE[E]$ pool at its own measured wall-clock cost, so both sides of the
ratio are charged like-for-like:
\begin{equation}
\FOM_{\rm wall}(\WE[I_\theta])\big/\FOM_{\rm wall}(\WE[E]) = 0.33,
\qquad 95\%\ \mathrm{CI}=[0.19,0.58],
\label{eq:fomwall}
\end{equation}
below parity and excluding $1$. A learned coordinate can therefore be dynamically more
informative while simultaneously being less efficient per unit of computer time. At the tested
scale, repeated neural inference is the dominant measured limitation on wall-clock efficiency;
caching, batching, or a cheaper architecture are therefore natural targets for a dedicated
follow-up efficiency study.

\subsection{Relation to neighboring rare-event methods}\label{sec:relation}
The present implementation is one member of a broader splitting and path-sampling family.
FFS and AMS advance trajectories through fixed interfaces or adaptively chosen levels
\citep{allen2009,valeriani2007,cerou2007}, and transition-path sampling works directly with
reactive-path ensembles, which frees it from a fixed scalar coordinate but targets different
estimators than the stationary tail probabilities studied here
\citep{chandler1978,dellago1998,bolhuis2002}. WE differs in maintaining a continuously weighted
population and resampling within bins, so that the bins influence variance but not the target
measure \citep{huber1996,zhang2010,zuckerman2017}. What the transport analogy adds is a
principled target for choosing those bins: CADIS couples an importance map to source biasing
and weight windows \citep{wagner1998,haghighat2003,wagnerfw2014}, and the present
implementation adopts the population-control half of that logic without altering the
Metropolis kernel.

Against the neural rare-event literature, the distinction is where the learned object is
inserted. Flow-based and autoregressive samplers learn global equilibrium structure
\citep{wu2019,noe2019,nicoli2020,gabrie2022}; neural committor and Feynman--Kac methods target
dynamical progress toward a specific event \citep{khoo2019,li2022,strahan2023}; and active or
self-consistent committor schemes concentrate training and sampling effort in
transition-relevant regions \citep{rotskoff2022,kang2024,trizio2025}. The framework of Kim and
Cai is closest in spirit, combining a learned event-specific importance function with biased
transitions and branching \citep{kimcai2026}: in the language of \S\ref{sec:adjoint}, their
importance function plays the role of $h$, their modified transitions the role of the tilt, and
their branching the role of splitting. All of these place the learned function inside the
dynamics. The present approach places it only in the binning instead, using the unbiased WE resampling
construction to preserve the target expectation while making coordinate quality measurable
without confounding it with likelihood-ratio error.

\subsection{Scope and reproducibility}\label{sec:scope}
The principal learned-coordinate experiment uses six two-dimensional EA disorder
realizations, one temperature ($T=2.6$), one lattice size ($L=16$), one threshold depth
(beyond$=4$), and a compact periodic convolutional network. Replication is across trained
networks, disorder realizations, and simulation replicas, not yet across lattice size or
temperature. A larger disorder sample, to tighten the disorder-level effect estimate and better
characterize how its magnitude varies with disorder, and extensions in $L$ and $T$, are natural
next steps. The three independent network initializations used for the principal within-disorder
test produced consistent discovery rates ($63.3\%$--$70.0\%$), while the broader replication
strategy treats training initialization, quenched disorder, and simulation replicas as distinct
sources of uncertainty.

\FloatBarrier

% ============================================================================
\section{Conclusions}\label{sec:conclusions}

Adjoint transport theory provides a useful organizing principle for rare-event sampling in
lattice spin systems. The committor supplies the event-specific importance function, the exact
Doob transform defines the zero-variance ideal, and weight-conserving population control
provides a practical way to exploit approximate importance without changing the target
measure; these elements were validated sequentially from a solvable birth--death process to
exact enumeration of the Ising model.

At system scale, WE and naive sampling resolve the $L=20$ Ising magnetization tail to
comparable depth at $R=20$. Direct sampling is more precise in the near tail, while WE holds a
modest, roughly two-fold figure-of-merit advantage at the single deepest resolved bin. The
benefit of population control is therefore rarity-dependent rather than uniform across the
tail.

In the frustrated EA model, the learned configuration-space importance coordinate produces the
main new result of this study. Pooled across six independent disorder realizations with one
trained network each, it reaches the deep target in $121/180$ replicas against $95/180$ for the
physical energy coordinate ($67.2\%$ versus $52.8\%$, Fisher exact $p=0.0071$), reproducing a
within-disorder result across three independently trained networks against a balanced baseline
($66.7\%$ versus $48.9\%$, $p=0.023$). A learned progress coordinate can extract dynamical
information not fully represented by a scalar physical observable, and the advantage is not
confined to one particular quenched-disorder sample: it is reproduced across the six disorder
realizations tested.

That advantage is presently one of reliability rather than efficiency: repeated network
inference is directly measured to cost four to five times $\WE[E]$'s wall-clock time, leaving
the learned coordinate at roughly one-third of parity in figure of merit ($0.33$, $95\%$ CI
$[0.19,0.58]$). At the tested scale, repeated neural inference is therefore the dominant measured
limitation on wall-clock efficiency.

% ============================================================================
\section*{Data and Code Availability}
The code, input data, trained network weights, quenched disorder realizations, per-replica
outputs, validation data, and wall-clock records required to reproduce the reported results have
been assembled into an archival package. A persistent public repository DOI will be inserted
before final submission. The package also records the software environment, configuration files,
and random seeds used for the production calculations so that the reported analyses can be
reproduced under the archived environment.

\section*{Acknowledgements}
The author gratefully acknowledges the support of Fahad Bin Sultan University, Tabuk, Saudi
Arabia, including access to the computational environment used for the numerical experiments.

\printcredits

% TODO before submission: insert an author-confirmed Declaration of competing interest if
% required by the target journal. No statement is inserted automatically here because it is an
% author declaration and must not be inferred.

\section*{Declaration of generative AI and AI-assisted technologies in the manuscript preparation process}
During the preparation of this work, the author used ChatGPT (OpenAI) and Claude (Anthropic)
to assist with manuscript organization, language refinement, critical review, and software
support. After using these tools, the author reviewed and edited the content as needed,
independently verified the scientific claims and numerical results, and takes full
responsibility for the content of the publication.

% ============================================================================
\begin{thebibliography}{99}
\bibitem{bolhuis2002} P. G. Bolhuis, D. Chandler, C. Dellago, P. L. Geissler, Transition path
sampling: throwing ropes over rough mountain passes, in the dark, Annu. Rev. Phys. Chem. 53 (2002)
291--318.
\bibitem{kahn1951} H. Kahn, T. E. Harris, Estimation of particle transmission by random sampling,
Natl. Bur. Stand. Appl. Math. Ser. 12 (1951) 27--30.
\bibitem{kalos2008} M. H. Kalos, P. A. Whitlock, Monte Carlo Methods, 2nd ed., Wiley-VCH,
Weinheim, 2008.
\bibitem{e2006} W. E, E. Vanden-Eijnden, Towards a theory of transition paths, J. Stat. Phys. 123
(2006) 503--523.
\bibitem{metzner2009} P. Metzner, C. Sch\"utte, E. Vanden-Eijnden, Transition path theory for
Markov jump processes, Multiscale Model. Simul. 7 (2009) 1192--1219.
\bibitem{vandeneijnden2010} E. Vanden-Eijnden, Transition path theory, in: M. Ferrario,
G. Ciccotti, K. Binder (Eds.), Computer Simulations in Condensed Matter Systems: From Materials to
Chemical Biology, Vol. 1, Springer, 2010, pp. 453--493.
\bibitem{lewins1965} J. Lewins, Importance: The Adjoint Function, Pergamon Press, Oxford, 1965.
\bibitem{spanier1969} J. Spanier, E. M. Gelbard, Monte Carlo Principles and Neutron Transport
Problems, Addison-Wesley, Reading, MA, 1969.
\bibitem{bell1970} G. I. Bell, S. Glasstone, Nuclear Reactor Theory, Van Nostrand Reinhold,
New York, 1970.
\bibitem{wagner1998} J. C. Wagner, A. Haghighat, Automated variance reduction of Monte Carlo
shielding calculations using the discrete ordinates adjoint function, Nucl. Sci. Eng. 128 (1998)
186--208.
\bibitem{haghighat2003} A. Haghighat, J. C. Wagner, Monte Carlo variance reduction with
deterministic importance functions, Prog. Nucl. Energy 42 (2003) 25--53.
\bibitem{doob1984} J. L. Doob, Classical Potential Theory and Its Probabilistic Counterpart,
Springer, New York, 1984.
\bibitem{touchette2009} H. Touchette, The large deviation approach to statistical mechanics,
Phys. Rep. 478 (2009) 1--69.
\bibitem{chetrite2015} R. Chetrite, H. Touchette, Nonequilibrium Markov processes conditioned on
large deviations, Ann. Henri Poincar\'e 16 (2015) 2005--2057.
\bibitem{wagnerfw2014} J. C. Wagner, D. E. Peplow, S. W. Mosher, FW-CADIS method for global and
regional variance reduction of Monte Carlo radiation transport calculations, Nucl. Sci. Eng. 176
(2014) 37--57.
\bibitem{huber1996} G. A. Huber, S. Kim, Weighted-ensemble Brownian dynamics simulations for
protein association reactions, Biophys. J. 70 (1996) 97--110.
\bibitem{zhang2010} B. W. Zhang, D. Jasnow, D. M. Zuckerman, The ``weighted ensemble'' path
sampling method is statistically exact for a broad class of stochastic processes and binning
procedures, J. Chem. Phys. 132 (2010) 054107.
\bibitem{bhatt2010} D. Bhatt, B. W. Zhang, D. M. Zuckerman, Steady-state simulations using
weighted ensemble path sampling, J. Chem. Phys. 133 (2010) 014110.
\bibitem{russo2022} J. D. Russo, et al., WESTPA 2.0: high-performance upgrades for weighted
ensemble simulations and analysis of longer-timescale applications, J. Chem. Theory Comput. 18
(2022) 638--649.
\bibitem{allen2009} R. J. Allen, C. Valeriani, P. R. ten Wolde, Forward flux sampling for rare
event simulations, J. Phys.: Condens. Matter 21 (2009) 463102.
\bibitem{valeriani2007} C. Valeriani, R. J. Allen, M. J. Morelli, D. Frenkel, P. R. ten Wolde,
Computing stationary distributions in equilibrium and nonequilibrium systems with forward flux
sampling, J. Chem. Phys. 127 (2007) 114109.
\bibitem{cerou2007} F. C\'erou, A. Guyader, Adaptive multilevel splitting for rare event analysis,
Stoch. Anal. Appl. 25 (2007) 417--443.
\bibitem{dellago1998} C. Dellago, P. G. Bolhuis, D. Chandler, Efficient transition path sampling:
application to Lennard-Jones cluster rearrangements, J. Chem. Phys. 108 (1998) 9236--9245.
\bibitem{wu2019} D. Wu, L. Wang, P. Zhang, Solving statistical mechanics using variational
autoregressive networks, Phys. Rev. Lett. 122 (2019) 080602.
\bibitem{noe2019} F. No\'e, S. Olsson, J. K\"ohler, H. Wu, Boltzmann generators: sampling
equilibrium states of many-body systems with deep learning, Science 365 (2019) eaaw1147.
\bibitem{nicoli2020} K. A. Nicoli, et al., Asymptotically unbiased estimation of physical
observables with neural samplers, Phys. Rev. E 101 (2020) 023304.
\bibitem{gabrie2022} M. Gabri\'e, G. M. Rotskoff, E. Vanden-Eijnden, Adaptive Monte Carlo augmented
with normalizing flows, Proc. Natl. Acad. Sci. U.S.A. 119 (2022) e2109420119.
\bibitem{khoo2019} Y. Khoo, J. Lu, L. Ying, Solving for high-dimensional committor functions using
artificial neural networks, Res. Math. Sci. 6 (2019) 1.
\bibitem{strahan2023} J. Strahan, J. Finkel, A. R. Dinner, J. Weare, Predicting rare events using
neural networks and short-trajectory data, J. Comput. Phys. 488 (2023) 112152.
\bibitem{hua2024} X. Hua, R. Ahmad, J. Blanchet, W. Cai, Accelerated sampling of rare events using
a neural network bias potential, arXiv:2401.06936 (2024).
\bibitem{kimcai2026} M. Kim, W. Cai, Accelerated Markov Chain Monte Carlo simulation via neural
network-driven importance sampling, arXiv:2602.12294 (2026).
\bibitem{zhao2025} Y. Zhao, Y. Tang, P. Zhang, Nonequilibrium statistical mechanics revealed by
Doob $h$ transform and variational autoregressive networks, Phys. Rev. E 111 (2025) 034120.
\bibitem{ising1925} E. Ising, Beitrag zur Theorie des Ferromagnetismus, Z. Phys. 31 (1925)
253--258.
\bibitem{onsager1944} L. Onsager, Crystal statistics. I. A two-dimensional model with an
order-disorder transition, Phys. Rev. 65 (1944) 117--149.
\bibitem{edwards1975} S. F. Edwards, P. W. Anderson, Theory of spin glasses, J. Phys. F: Met.
Phys. 5 (1975) 965--974.
\bibitem{binder1986} K. Binder, A. P. Young, Spin glasses: experimental facts, theoretical
concepts, and open questions, Rev. Mod. Phys. 58 (1986) 801--976.
\bibitem{metropolis1953} N. Metropolis, A. W. Rosenbluth, M. N. Rosenbluth, A. H. Teller,
E. Teller, Equation of state calculations by fast computing machines, J. Chem. Phys. 21 (1953)
1087--1092.
\bibitem{zuckerman2017} D. M. Zuckerman, L. T. Chong, Weighted ensemble simulation: review of
methodology, applications, and software, Annu. Rev. Biophys. 46 (2017) 43--57.
\bibitem{li2022} H. Li, Y. Khoo, Y. Ren, L. Ying, A semigroup method for high dimensional committor
functions based on neural network, Proc. Mach. Learn. Res. 145 (2022) 598--618.
\bibitem{rotskoff2022} G. M. Rotskoff, A. R. Mitchell, E. Vanden-Eijnden, Active importance
sampling for variational objectives dominated by rare events: consequences for optimization and
generalization, Proc. Mach. Learn. Res. 145 (2022) 757--780.
\bibitem{kang2024} P. Kang, E. Trizio, M. Parrinello, Computing the committor with the committor to
study the transition state ensemble, Nat. Comput. Sci. 4 (2024) 451--460.
\bibitem{trizio2025} E. Trizio, P. Kang, M. Parrinello, Everything everywhere all at once: a
probability-based enhanced sampling approach to rare events, Nat. Comput. Sci. 5 (2025)
582--591.
\bibitem{mantel1959} N. Mantel, W. Haenszel, Statistical aspects of the analysis of data from
retrospective studies of disease, J. Natl. Cancer Inst. 22 (1959) 719--748.
\bibitem{fisher1922} R. A. Fisher, On the interpretation of $\chi^2$ from contingency tables,
and the calculation of $P$, J. R. Stat. Soc. 85 (1922) 87--94.
\bibitem{efron1979} B. Efron, Bootstrap methods: another look at the jackknife, Ann. Statist.
7 (1979) 1--26.
\bibitem{chandler1978} D. Chandler, Statistical mechanics of isomerization dynamics in liquids and
the transition state approximation, J. Chem. Phys. 68 (1978) 2959--2970.
\end{thebibliography}

\end{document}
